\documentclass[../main]{subfiles}
\begin{document}

\chapter{結論}
\label{cp:conclusion}
\thispagestyle{empty}
\minitoc
\newpage

\section{結論}
\label{sec:conc_conclusion}
本研究では，車椅子卓球競技クラス4の選手を対象に，打球の正確性を決定づける生体力学的要因を解明することを目的とし，計測システムの構築から決定論的モデルに基づく階層的な要因解析までを一貫して行った．
まず，本研究の貢献として，パフォーマンスと身体運動の因果関係を定量的に紐づけるための統合的な計測環境を実現したことが挙げられる．
本研究で構築したボールの位置，身体動作，および表面筋電図の完全同期システムは，各試行における着弾誤差を，その瞬間のラケット挙動，背後にある関節運動，そして駆動源である筋活動と1対1で対応させ，打球精度低下の原因を物理的かつ生理学的に追跡することを可能にした．
このシステムを用いて取得した多次元の時系列データに対して，機械学習および統計学的手法を用いた解析を行った結果，以下の知見が得られた．


まず，パフォーマンスを直接決定するラケット挙動の特定である．
本研究では，ラケットの運動学的パラメータを入力とした打球精度の分類モデルであるk-近傍法（KNN）を構築し，Permutation Feature Importance（PFI）を用いた特徴量の重要度算出を行った．この解析により，数ある変数の中でもラケット面角度，水平打ち出し角，および長軸方向の打点位置の3変数が，打球精度に対し影響を持つことが定量的に特定された．
とりわけ，打点位置の分布解析においては，精度の高い打球がラケットの先端側で捉えられているのに対し，打球精度の低い試行では，意図せずラケットの根本で打球してしまう打球の詰まりが確認された．これは，下肢による位置調整が困難なクラス4の車椅子卓球選手にとって，身体とボールとの間に適切な距離を確保できているかどうかが，精度の良し悪しを決定づける要因となることを示している．

次に，上記のエラー挙動を生じさせる身体動作メカニズムの解明を行うために，関節運動データに対するステップワイズ法を用いた重回帰分析，および筋活動量を用いた群間比較検定を行った．
その結果．打球ミスが単発的なフォームの崩れではなく，以下の要因が複合的に作用した結果であることを明らかにした．

まず，運動学的な観点からは，前の打球からの戻りが遅れることによる準備の不足が，打球時に詰まりを引き起こしていることが分かった．
解析の結果，準備が間に合わないまま遠くのボールを打とうとすると，選手は焦って上体を大きく倒し込んでしまう傾向が見られた．
しかし，統計データはこの動作が逆効果であることを示している．
無理に体を倒すことで，かえって肩の位置がボールに近づきすぎてしまい，結果として腕を伸び伸びと振るための空間を失ってしまうのである．

次に，この窮屈な状態は筋活動の面から見ても悪循環を生んでいた．
本来であれば，身体の中心から腕先へと順序よく力が伝わる運動連鎖が生じるはずであるが，打球精度が低下する試行においてはこの流れが途切れていた．スムーズに力が伝わらないため，選手は無意識に全身に力を入れて固めたり，手先だけでラケットを操作しようとしたりする．
これが筋電図で確認された過剰な力みの正体であり，この代償動作がスイング軌道を不安定にさせていたと結論付けられる．

つまり，パフォーマンス低下の原因は，打ち方そのものではなく，打つ前の準備の質にある．
打球精度を高めるために必要な要素として，インパクトの瞬間のフォームを修正に加えて，前のボールを打った直後に素早く戻り，次のボールに対して十分な時間的余裕と空間的余裕を確保するという，打つ前の準備のプロセスを改善することが，結果として理想的なスイングを生み出す鍵となる．
これは，障害特性に応じた指導を行う上で，静的なスイングの形だけでなく，一連の動作の流れとしてのリカバリー能力や運動連鎖のタイミングを重視すべきであることを示す新たな知見である．

\clearpage

\section{今後の展望}
\label{sec:conc_future}
本研究の今後の展望として，以下の二点が挙げられる．

第1に，データセットの拡張および健常者データとの比較による競技特性の明確化である．
本研究では単一のクラス4選手を対象としたが，構築したモデルの汎用性と信頼性を高めるためには，より多くの試行数を確保し，外れ値の影響を小さくしていくことで機械学習モデル等の解析精度を向上させる必要がある．
さらに重要な点は，健常者卓球選手のデータとの比較対照である．本研究で明らかになったラケット先端打撃の優位性や上部体幹主導の動作といった特徴が，個人の癖によるものなのか，それともクラス4という障害特性に起因する適応戦略なのかを厳密に区別する必要がある．
健常者の標準的なバイオメカニクスモデルと本被験者のモデルを比較することで，パラスポーツ特有の身体操作戦略をより際立たせ，障害を持った選手の運動戦略の理解をより深めることができると考えられる．

第2に，身体部位間の連動性に関するより詳細な解析である．
本研究では，個々の関節速度やタイミングの遅れに着目したが，今後は各身体セグメントがどのように協調してエネルギーを伝達しているかという連動性の定量化が課題となる．
具体的には，体幹から肩，肘，手首へと至るエネルギーフローの効率性や，関節間の位相差等を特徴量とした解析をすることで，なぜ戻り動作の遅れが末端のミスに直結するのかという力学的メカニズムをより深く解明できる．
特に，運動連鎖が破綻しやすい車椅子選手において，どの関節がどのタイミングで代償動作を行っているかを詳細に可視化することは，効果的なトレーニング介入を行う上で不可欠な要素となると考えられる．
\clearpage

\end{document}
